\import{set/suc.tex}
\import{function.tex}

\section{Natural numbers}

Inductive sets, the natural numbers, and the fact that $\emptyset$ belongs to
$\naturals$ are supplied directly by the standard prelude.  Natural addition
is the function represented by the least relation containing the zero cases
and closed under simultaneous successor in its second argument and value.

\begin{definition}\label{num_addition_bound}
    $\numAdditionBound =
        (\naturals\times\naturals)\times\naturals$.
\end{definition}

\begin{definition}\label{num_addition_closed_relations}
    $\numAdditionClosedRelations =
        \{R\in\pow{\numAdditionBound}\mid
            \text{for all $a\in\naturals$ we have
                $((a,\emptyset),a)\in R$ and
            for all $x,y,z$ such that $((x,y),z)\in R$ we have
                $((x,\suc{y}),\suc{z})\in R$}\}$.
\end{definition}

\begin{proposition}\label{num_addition_bound_closed}
    $\numAdditionBound\in\numAdditionClosedRelations$.
\end{proposition}
\begin{proof}
    $\numAdditionBound\in\pow{\numAdditionBound}$
        by \obv{powerset_intro,subseteq_refl}.
    For all $x\in\naturals$ we have
        $((x,\emptyset),x)\in\numAdditionBound$
        by \obv{num_addition_bound,times_tuple_intro,omega_empty}.
    For all $x,y,z$ such that
        $((x,y),z)\in\numAdditionBound$ we have
        $((x,\suc{y}),\suc{z})\in\numAdditionBound$
        by \auto{num_addition_bound,times_tuple_elim,times_tuple_intro,omega_successor}.
    Follows by \obv{num_addition_closed_relations}.
\end{proof}

\begin{proposition}\label{num_addition_closed_relations_inhabited}
    $\numAdditionClosedRelations$ is inhabited.
\end{proposition}
\begin{proof}
    Follows by \obv{num_addition_bound_closed}.
\end{proof}

\begin{definition}\label{num_addition_graph}
    $\numAdditionGraph =
        \inters{\numAdditionClosedRelations}$.
\end{definition}

\begin{proposition}\label{num_addition_graph_minimal}
    Let $R\in\numAdditionClosedRelations$.
    Then $\numAdditionGraph\subseteq R$.
\end{proposition}
\begin{proof}
    Follows by \auto{num_addition_graph,inters_subseteq_elem}.
\end{proof}

\begin{proposition}\label{num_addition_graph_bound}
    $\numAdditionGraph\subseteq\numAdditionBound$.
\end{proposition}
\begin{proof}
    Follows by \obv{num_addition_graph_minimal,num_addition_bound_closed}.
\end{proof}

\begin{proposition}\label{num_addition_graph_zero_intro}
    Let $x\in\naturals$.
    Then $((x,\emptyset),x)\in\numAdditionGraph$.
\end{proposition}
\begin{proof}
    For all $R\in\numAdditionClosedRelations$ we have
        $((x,\emptyset),x)\in R$
        by \obv{num_addition_closed_relations}.
    Follows by \obv{num_addition_graph,inters_intro,num_addition_closed_relations_inhabited}.
\end{proof}

\begin{proposition}\label{num_addition_graph_successor_intro}
    Suppose $((x,y),z)\in\numAdditionGraph$.
    Then $((x,\suc{y}),\suc{z})\in\numAdditionGraph$.
\end{proposition}
\begin{proof}
    For all $R\in\numAdditionClosedRelations$ we have
        $((x,y),z)\in R$ by \auto{num_addition_graph_minimal,subseteq}.
    For all $R\in\numAdditionClosedRelations$ we have
        $((x,\suc{y}),\suc{z})\in R$
        by \auto{num_addition_closed_relations}.
    Follows by \obv{num_addition_graph,inters_intro,num_addition_closed_relations_inhabited}.
\end{proof}

\begin{proposition}\label{num_addition_graph_closed}
    $\numAdditionGraph\in\numAdditionClosedRelations$.
\end{proposition}
\begin{proof}
    $\numAdditionGraph\in\pow{\numAdditionBound}$
        by \obv{powerset_intro,num_addition_graph_bound}.
    Follows by \obv{num_addition_closed_relations,num_addition_graph_zero_intro,num_addition_graph_successor_intro}.
\end{proof}

\begin{definition}\label{num_addition_bases}
    $\numAdditionBases =
        \{w\in\numAdditionBound\mid
            \text{there exists $x\in\naturals$ such that
                $w=((x,\emptyset),x)$}\}$.
\end{definition}

\begin{definition}\label{num_addition_steps}
    $\numAdditionSteps =
        \{w\in\numAdditionBound\mid
            \text{there exist $x,y,z$ such that
                $((x,y),z)\in\numAdditionGraph$ and
                $w=((x,\suc{y}),\suc{z})$}\}$.
\end{definition}

\begin{definition}\label{num_addition_generated}
    $\numAdditionGenerated =
        \numAdditionBases\union\numAdditionSteps$.
\end{definition}

\begin{proposition}\label{num_addition_generated_bound}
    $\numAdditionGenerated\subseteq\numAdditionBound$.
\end{proposition}
\begin{proof}
    Follows by \auto{num_addition_generated,num_addition_bases,num_addition_steps,subseteq,union_iff}.
\end{proof}

\begin{proposition}\label{num_addition_generated_subset_graph}
    $\numAdditionGenerated\subseteq\numAdditionGraph$.
\end{proposition}
\begin{proof}
    Follows by \auto{num_addition_generated,num_addition_bases,num_addition_steps,subseteq,union_iff,num_addition_graph_zero_intro,num_addition_graph_successor_intro}.
\end{proof}

\begin{proposition}\label{num_addition_generated_closed}
    $\numAdditionGenerated\in\numAdditionClosedRelations$.
\end{proposition}
\begin{proof}
    $\numAdditionGenerated\in\pow{\numAdditionBound}$
        by \obv{powerset_intro,num_addition_generated_bound}.
    For all $x\in\naturals$ we have
        $((x,\emptyset),x)\in\numAdditionBound$
        by \obv{num_addition_bound,times_tuple_intro,omega_empty}.
    For all $x\in\naturals$ we have
        $((x,\emptyset),x)\in\numAdditionBases$
        by \obv{num_addition_bases}.
    For all $x\in\naturals$ we have
        $((x,\emptyset),x)\in\numAdditionGenerated$
        by \auto{num_addition_generated,union_iff}.
    For all $x,y,z$ such that
        $((x,y),z)\in\numAdditionGenerated$ we have
        $((x,y),z)\in\numAdditionGraph$
        by \obv{num_addition_generated_subset_graph,subseteq}.
    For all $x,y,z$ such that
        $((x,y),z)\in\numAdditionGenerated$ we have
        $((x,y),z)\in\numAdditionBound$
        by \obv{num_addition_graph_bound,subseteq}.
    For all $x,y,z$ such that
        $((x,y),z)\in\numAdditionGenerated$ we have
        $((x,\suc{y}),\suc{z})\in\numAdditionBound$
        by \auto{num_addition_bound,times_tuple_elim,times_tuple_intro,omega_successor}.
    For all $x,y,z$ such that
        $((x,y),z)\in\numAdditionGenerated$ we have
        $((x,\suc{y}),\suc{z})\in\numAdditionSteps$
        by \auto{num_addition_steps}.
    For all $x,y,z$ such that
        $((x,y),z)\in\numAdditionGenerated$ we have
        $((x,\suc{y}),\suc{z})\in\numAdditionGenerated$
        by \auto{num_addition_generated,union_iff}.
    Follows by \obv{num_addition_closed_relations}.
\end{proof}

\begin{proposition}\label{num_addition_graph_cases}
    Suppose $w\in\numAdditionGraph$.
    Then $w\in\numAdditionBases$ or $w\in\numAdditionSteps$.
\end{proposition}
\begin{proof}
    $\numAdditionGraph\subseteq\numAdditionGenerated$
        by \obv{num_addition_graph_minimal,num_addition_generated_closed}.
    Follows by \auto{subseteq,num_addition_generated,union_iff}.
\end{proof}

\begin{proposition}\label{num_addition_graph_zero_elim}
    Suppose $((x,\emptyset),z)\in\numAdditionGraph$.
    Then $z=x$.
\end{proposition}
\begin{proof}
    We have $((x,\emptyset),z)\in\numAdditionBases$ or
        $((x,\emptyset),z)\in\numAdditionSteps$
        by \obv{num_addition_graph_cases}.
    Follows by \auto{num_addition_bases,num_addition_steps,pair_eq_iff,suc_neq_emptyset}.
\end{proof}

\begin{proposition}\label{num_addition_graph_successor_elim}
    Suppose $((x,\suc{y}),z)\in\numAdditionGraph$.
    Then there exists $w$ such that
        $((x,y),w)\in\numAdditionGraph$ and $z=\suc{w}$.
\end{proposition}
\begin{proof}
    We have $((x,\suc{y}),z)\in\numAdditionBases$ or
        $((x,\suc{y}),z)\in\numAdditionSteps$
        by \obv{num_addition_graph_cases}.
    Follows by \auto{num_addition_bases,num_addition_steps,pair_eq_iff,suc_neq_emptyset,suc_injective}.
\end{proof}

\begin{proposition}\label{num_addition_graph_total}
    Let $x,y\in\naturals$.
    Then there exists $z\in\naturals$ such that
        $((x,y),z)\in\numAdditionGraph$.
\end{proposition}
\begin{proof}
    Let $D=\{u\in\naturals\mid
        \text{there exists $v\in\naturals$ such that
            $((x,u),v)\in\numAdditionGraph$}\}$.
    $\emptyset\in D$
        by \obv{omega_empty,num_addition_graph_zero_intro}.
    For all $u\in D$ we have $\suc{u}\in D$
        by \auto{omega_successor,num_addition_graph_successor_intro}.
    Thus $D$ is an inductive set by \obv{inductive_set}.
    Hence every element of $\naturals$ is an element of $D$
        by \obv{omega_minimal}.
    Thus $y\in D$.
\end{proof}

\begin{proposition}\label{num_addition_graph_unique}
    Let $x,y\in\naturals$.
    Suppose $((x,y),z),((x,y),z')\in\numAdditionGraph$.
    Then $z=z'$.
\end{proposition}
\begin{proof}
    Let $U=\{u\in\naturals\mid
        \text{for all $v,w$ such that
            $((x,u),v),((x,u),w)\in\numAdditionGraph$
            we have $v=w$}\}$.
    $\emptyset\in U$
        by \obv{omega_empty,num_addition_graph_zero_elim}.
    For all $u\in U$ we have $\suc{u}\in U$
        by \auto{omega_successor,num_addition_graph_successor_elim,suc_injective}.
    Thus $U$ is an inductive set by \obv{inductive_set}.
    Hence every element of $\naturals$ is an element of $U$
        by \obv{omega_minimal}.
    Thus $y\in U$.
\end{proof}

\begin{proposition}\label{num_addition_graph_relation}
    $\numAdditionGraph$ is a relation.
\end{proposition}
\begin{proof}
    Follows by \auto{relation,num_addition_graph_bound,num_addition_bound,subseteq,times_elem_is_tuple}.
\end{proof}

\begin{proposition}\label{num_addition_graph_rightunique}
    $\numAdditionGraph$ is right-unique.
\end{proposition}
\begin{proof}
    It suffices to show that for all $p,z,z'$ such that
        $(p,z),(p,z')\in\numAdditionGraph$ we have $z=z'$
        by \obv{rightunique}.
    Fix $p,z,z'$.
    Assume $(p,z),(p,z')\in\numAdditionGraph$.
    We have $(p,z)\in\numAdditionBound$
        by \obv{num_addition_graph_bound,subseteq}.
    Thus $p\in\naturals\times\naturals$
        by \auto{num_addition_bound,times_tuple_elim}.
    Take $x,y$ such that $p=(x,y)$ and $x,y\in\naturals$
        by \auto{times_elem_is_tuple}.
    Then $((x,y),z),((x,y),z')\in\numAdditionGraph$.
    Follows by \obv{num_addition_graph_unique}.
\end{proof}

\begin{proposition}\label{num_addition_graph_function}
    $\numAdditionGraph$ is a function.
\end{proposition}
\begin{proof}
    Follows by \obv{num_addition_graph_relation,num_addition_graph_rightunique}.
\end{proof}

\begin{proposition}\label{num_addition_graph_domain}
    $\dom{\numAdditionGraph}=\naturals\times\naturals$.
\end{proposition}
\begin{proof}
    Every element of $\dom{\numAdditionGraph}$ is an element of
        $\naturals\times\naturals$
        by \auto{dom_iff,num_addition_graph_bound,num_addition_bound,subseteq,times_tuple_elim}.
    Every element of $\naturals\times\naturals$ is an element of
        $\dom{\numAdditionGraph}$
        by \auto{times_elem_is_tuple,num_addition_graph_total,dom_iff}.
    Follows by set extensionality.
\end{proof}

\begin{proposition}\label{num_addition_graph_funs}
    $\numAdditionGraph\in
        \funs{\naturals\times\naturals}{\naturals}$.
\end{proposition}
\begin{proof}
    For all $p\in\dom{\numAdditionGraph}$ we have
        $(p,\apply{\numAdditionGraph}{p})\in\numAdditionGraph$
        by \obv{num_addition_graph_function,function_apply_elim}.
    For all $p\in\dom{\numAdditionGraph}$ we have
        $(p,\apply{\numAdditionGraph}{p})\in\numAdditionBound$
        by \obv{num_addition_graph_bound,subseteq}.
    For all $p\in\dom{\numAdditionGraph}$ we have
        $\apply{\numAdditionGraph}{p}\in\naturals$
        by \auto{num_addition_bound,times_tuple_elim}.
    Thus $\numAdditionGraph$ is a function to $\naturals$
        by \obv{num_addition_graph_function}.
    Follows by \obv{funs_intro,num_addition_graph_domain}.
\end{proof}

\begin{abbreviation}\label{num_zero_is_emptyset}
    $\zero = \emptyset$.
\end{abbreviation}

\begin{definition}\label{num_addition_is_set}
    $x+y=\apply{\numAdditionGraph}{(x,y)}$.
\end{definition}

\begin{classification}\label{num_addition_on_naturals}
    Suppose $x,y\in\naturals$.
    Then $x+y\in\naturals$.
\end{classification}
\begin{proof}
    $(x,y)\in\naturals\times\naturals$
        by \obv{times_tuple_intro}.
    $\apply{\numAdditionGraph}{(x,y)}\in\naturals$
        by \obv{num_addition_graph_funs,funs_type_apply}.
    Follows by \obv{num_addition_is_set}.
\end{proof}

\begin{proposition}\label{num_addition_zero_right}
    For all $x\in\naturals$ we have $x+\emptyset=x$.
\end{proposition}
\begin{proof}
    Follows by \obv{num_addition_is_set,num_addition_graph_function,num_addition_graph_zero_intro,function_apply_intro}.
\end{proof}

\begin{proposition}\label{num_addition_successor_right}
    Let $x,y\in\naturals$.
    Then $x+\suc{y}=\suc{x+y}$.
\end{proposition}
\begin{proof}
    $(x,y)\in\dom{\numAdditionGraph}$
        by \obv{num_addition_graph_domain,times_tuple_intro}.
    $((x,y),\apply{\numAdditionGraph}{(x,y)})
        \in\numAdditionGraph$
        by \obv{num_addition_graph_function,function_apply_elim}.
    $((x,\suc{y}),\suc{\apply{\numAdditionGraph}{(x,y)}})
        \in\numAdditionGraph$
        by \obv{num_addition_graph_successor_intro}.
    $\apply{\numAdditionGraph}{(x,\suc{y})}
        =\suc{\apply{\numAdditionGraph}{(x,y)}}$
        by \obv{num_addition_graph_function,function_apply_intro}.
    Follows by \auto{num_addition_is_set}.
\end{proof}

\begin{proposition}\label{num_addition_zero_left}
    For all $x\in\naturals$ we have $\emptyset+x=x$.
\end{proposition}
\begin{proof}
    Let $Z=\{u\in\naturals\mid \text{$\emptyset+u=u$}\}$.
    $\emptyset\in Z$ by \obv{omega_empty,num_addition_zero_right}.
    For all $u\in Z$ we have $\suc{u}\in Z$
        by \auto{omega_successor,num_addition_successor_right}.
    Thus $Z$ is an inductive set by \obv{inductive_set}.
    Every element of $\naturals$ is an element of $Z$
        by \obv{omega_minimal}.
    Follows by \auto.
\end{proof}

\begin{proposition}\label{num_addition_successor_left}
    Let $x,y\in\naturals$.
    Then $\suc{x}+y=\suc{x+y}$.
\end{proposition}
\begin{proof}
    Let $S=\{u\in\naturals\mid
        \text{$\suc{x}+u=\suc{x+u}$}\}$.
    $\emptyset\in S$
        by \auto{omega_empty,omega_successor,num_addition_zero_right}.
    For all $u\in S$ we have $\suc{u}\in S$
        by \auto{omega_successor,num_addition_successor_right}.
    Thus $S$ is an inductive set by \obv{inductive_set}.
    Every element of $\naturals$ is an element of $S$
        by \obv{omega_minimal}.
    Follows by \auto.
\end{proof}

\begin{proposition}\label{num_addition_axiom_1}
    For all $x\in\naturals$ we have
        $x+\zero=\zero+x=x$.
\end{proposition}
\begin{proof}
    Follows by \obv{num_addition_zero_right,num_addition_zero_left}.
\end{proof}

\begin{proposition}\label{num_addition_axiom_2}
    Let $x,y\in\naturals$.
    Then $x+\suc{y}=\suc{x}+y=\suc{x+y}$.
\end{proposition}
\begin{proof}
    Follows by \obv{num_addition_successor_right,num_addition_successor_left}.
\end{proof}

\begin{theorem}\label{natural_addition_commutative}
    Let $m,n$ be natural numbers.
    Then $m+n=n+m$.
\end{theorem}
\begin{proof}
    Let $C=\{k\in\naturals\mid \text{$m+k=k+m$}\}$.
    $\emptyset\in C$
        by \obv{omega_empty,num_addition_zero_right,num_addition_zero_left}.
    Show that for every $k\in C$ we have $\suc{k}\in C$.
    \begin{subproof}
        Fix $k\in C$.
        Then $k$ is a natural number and $m+k=k+m$.
        $m+\suc{k}=\suc{m+k}$
            by \obv{num_addition_successor_right}.
        $\suc{k}+m=\suc{k+m}$
            by \obv{num_addition_successor_left}.
        $\suc{k}$ is a natural number by \obv{omega_successor}.
    \end{subproof}
    Hence $C$ is an inductive set by \obv{inductive_set}.
    Hence every natural number is an element of $C$
        by \obv{omega_minimal}.
    Thus $n\in C$.
\end{proof}

\begin{theorem}\label{natural_addition_associative}
    Let $l,m,n$ be natural numbers.
    Then $(l+m)+n=l+(m+n)$.
\end{theorem}
\begin{proof}
    $l+m$ is a natural number by \obv{num_addition_on_naturals}.
    Let $C=\{k\in\naturals\mid
        \text{$(l+m)+k=l+(m+k)$}\}$.
    $\emptyset\in C$
        by \obv{omega_empty,num_addition_zero_right}.
    Show that for every $k\in C$ we have $\suc{k}\in C$.
    \begin{subproof}
        Fix $k\in C$.
        Then $k$ is a natural number and
            $(l+m)+k=l+(m+k)$.
        $m+k$ is a natural number by \obv{num_addition_on_naturals}.
        $(l+m)+\suc{k}=\suc{(l+m)+k}$ and
            $l+(m+\suc{k})=\suc{l+(m+k)}$
            by \auto{num_addition_successor_right}.
        $\suc{k}$ is a natural number by \obv{omega_successor}.
        Follows by \obv{num_addition_successor_right}.
    \end{subproof}
    Hence $C$ is an inductive set by \obv{inductive_set}.
    Hence every natural number is an element of $C$
        by \obv{omega_minimal}.
    Thus $n\in C$.
\end{proof}

\begin{theorem}\label{natural_addition_cancel_right}
    Let $x,y,z$ be natural numbers.
    Suppose $y+x=z+x$.
    Then $y=z$.
\end{theorem}
\begin{proof}
    Let $C=\{k\in\naturals\mid
        \text{if $y+k=z+k$, then $y=z$}\}$.
    Show that $\emptyset\in C$.
    \begin{subproof}
        $\emptyset$ is a natural number by \obv{omega_empty}.
        Show that if $y+\emptyset=z+\emptyset$, then $y=z$.
        \begin{subproof}
            Assume $y+\emptyset=z+\emptyset$.
            Follows by \obv{num_addition_zero_right}.
        \end{subproof}
    \end{subproof}
    Show that for every $k\in C$ we have $\suc{k}\in C$.
    \begin{subproof}
        Fix $k\in C$.
        If $y+k=z+k$, then $y=z$.
        $\suc{k}$ is a natural number by \obv{omega_successor}.
        Show that if $y+\suc{k}=z+\suc{k}$, then $y=z$.
        \begin{subproof}
            Assume $y+\suc{k}=z+\suc{k}$.
            $\suc{y+k}=\suc{z+k}$
                by \obv{num_addition_successor_right}.
            Hence $y+k=z+k$ by \obv{suc_injective}.
        \end{subproof}
    \end{subproof}
    Hence $C$ is an inductive set by \obv{inductive_set}.
    Hence every natural number is an element of $C$
        by \obv{omega_minimal}.
    Thus $x\in C$.
\end{proof}

\begin{corollary}\label{natural_addition_cancel_left}
    Let $x,y,z$ be natural numbers.
    Suppose $x+y=x+z$.
    Then $y=z$.
\end{corollary}
\begin{proof}
    $y+x=z+x$ by \obv{natural_addition_commutative}.
    Follows by \obv{natural_addition_cancel_right}.
\end{proof}

\begin{theorem}\label{natural_addition_left_subseteq_sum}
    Let $m,k$ be natural numbers.
    Then $m\subseteq m+k$.
\end{theorem}
\begin{proof}
    Let $C=\{r\in\naturals\mid \text{$m\subseteq m+r$}\}$.
    $\emptyset\in C$ by
        \obv{omega_empty,num_addition_zero_right,subseteq_refl}.
    Show that for every $r\in C$ we have $\suc{r}\in C$.
    \begin{subproof}
        Fix $r\in C$.
        Then $r$ is a natural number and $m\subseteq m+r$.
        $\suc{r}$ is a natural number by \obv{omega_successor}.
        $m+r\subseteq\suc{m+r}$ by \obv{subseteq_self_suc_intro}.
        Hence $m\subseteq\suc{m+r}$ by \obv{subseteq_transitive}.
        $m+\suc{r}=\suc{m+r}$
            by \obv{num_addition_successor_right}.
        Thus $m\subseteq m+\suc{r}$.
    \end{subproof}
    Hence $C$ is an inductive set by \obv{inductive_set}.
    Hence every natural number is an element of $C$
        by \obv{omega_minimal}.
    Thus $k\in C$.
\end{proof}

\begin{lemma}\label{num_naturals_is_equal_to_two_times_naturals}
    $\{x+y \mid x \in \naturals, y \in \naturals \} = \naturals$.
\end{lemma}
\begin{proof}
    We have for all $z\in \{x+y \mid x \in \naturals, y \in \naturals \}$
        we have $z\in\naturals$ by \auto{num_addition_on_naturals}.
    We have for all $z\in\naturals$
        we have $z\in \{x+y \mid x \in \naturals, y \in \naturals \}$
        by \auto{num_addition_axiom_1,omega_empty}.
    Follows by set extensionality.
\end{proof}

\begin{theorem}\label{natural_zero_or_successor}
    Let $n$ be a natural number.
    Then $n=\emptyset$ or there exists a natural number $m$ such that
        $n=\suc{m}$.
\end{theorem}
\begin{proof}
    Let $C=\{k\in\naturals\mid
        \text{$k=\emptyset$ or there exists $j\in\naturals$ such that
            $k=\suc{j}$}\}$.
    $\emptyset\in C$ by \obv{omega_empty}.
    For every $k\in C$ we have $\suc{k}\in C$
        by \obv{omega_successor}.
    Hence $C$ is an inductive set by \obv{inductive_set}.
    Hence every natural number is an element of $C$
        by \obv{omega_minimal}.
    Thus $n\in C$.
\end{proof}
